\respcap{cap:dijkstra}
\resp{dj-bfs} Por indução nos vértices removidos: com pesos 1, os rótulos dos vértices na fila de prioridades assumem, em cada momento, no máximo dois valores consecutivos, $d$ e $d+1$, e os de valor $d$ saem antes, como na fila da busca em largura. Quando $v$ sai com $L(v)=d$, cada vizinho $w$ ainda não alcançado recebe $L(w)=d+1$ e $pai(w)=v$, exatamente como na busca em largura; vizinhos já alcançados têm $L(w)\le d+1$ e não mudam. Logo $L(w)=nivel(w)=\operatorname{dist}(raiz,w)$, e a árvore é de largura (com os empates resolvidos na mesma ordem).

\resp{dj-nd} Transforme cada aresta $xy$ de peso $p$ em duas arestas direcionadas, $x\to y$ e $y\to x$, ambas de peso $p$ (e use $\infty$ onde não há aresta), e execute o algoritmo no digrafo resultante.

\resp{dj-slide} Ver o Exemplo da §\ref{sec:dijkstra}: são removidos $a\,(0)$, $b\,(1)$, $c\,(2)$, $d\,(3)$ e $e\,(5)$, e a árvore é $a\to b\to c\to d\to e$.

\resp{dj-sim} São removidos $s\,(0)$, $y\,(5)$, $z\,(7)$, $t\,(8)$ e $x\,(9)$. Pais: $y\leftarrow s$, $z\leftarrow y$, $t\leftarrow y$ e $x\leftarrow t$. Por exemplo, o caminho mínimo até $x$ é $s\to y\to t\to x$, com custo $5+3+1$.

\resp{dj-neg} Com $x\to y$ (1), $x\to z$ (2) e $z\to y$ ($-2$), o algoritmo remove $x$, depois $y$ com $L(y)=1$, e o fixa; mas $x\to z\to y$ custa 0. Somar uma constante $c$ a todos os pesos \textbf{não} resolve: um caminho com $k$ arestas ganha $kc$, o que favorece os caminhos com menos arestas. No exemplo, com $c=2$: $x\to y$ passa a valer 3 e $x\to z\to y$ passa a valer $4+0=4$; o caminho ótimo muda.

\resp{dj-conf} A probabilidade de um caminho é o produto $\prod p(e)$; maximizá-la equivale a minimizar $\sum-\log p(e)$. Como $p(e)\in(0,1]$, os pesos $-\log p(e)$ são não negativos, e o algoritmo de Dijkstra se aplica. (Alternativa: um Dijkstra que mantém o maior produto e remove sempre o maior rótulo, o que funciona porque os fatores são no máximo 1.)

\resp{dj-fontes} Crie uma \emph{superfonte} $s^*$, com arestas de peso 0 para cada $s_i$, e execute o algoritmo a partir de $s^*$; ou, de forma equivalente, comece com $L(s_i)=0$ para todas as fontes, todas inseridas na fila de prioridades. A árvore resultante também indica qual fonte é a mais próxima de cada vértice.

\resp{dj-escolha} No primeiro grafo, esparso, use heap binário: $(n+m)\log n\approx4\cdot10^5\cdot17\approx7\cdot10^6$ operações, contra $n^2=10^{10}$ com vetor. No segundo, denso ($m\approx n^2/2$), use vetor: $n^2=10^8$, contra $(n+m)\log n\approx5\cdot10^7\cdot14=7\cdot10^8$ com heap.

\resp{dj-und} São removidos $1\,(0)$, $2\,(7)$, $3\,(9)$, $6\,(11)$, $4\,(20)$ e $5\,(20)$. Pais: $2\leftarrow1$, $3\leftarrow1$, $6\leftarrow3$, $4\leftarrow3$ e $5\leftarrow6$. Por exemplo, o caminho mínimo até 5 é $1,3,6,5$, com custo $9+2+9=20$, mais barato que $1\to6\to5$, que custa 23.

\resp{dj-zero} Sim. A prova de correção só usa que o trecho de $y$ até $u$ num caminho mínimo custa pelo menos 0, o que continua valendo com pesos nulos.

\resp{dj-alvo} Pode parar assim que $t$ for removido da fila de prioridades: pelo teorema de correção, nesse momento $L(t)=\operatorname{dist}(s,t)$, e as iterações seguintes só fixam vértices mais distantes, sem alterar $L(t)$.

\resp{dj-vert} Divida cada vértice $v$ em dois, $v_{in}$ e $v_{out}$, ligados por uma aresta $v_{in}\to v_{out}$ de peso $c(v)$; cada aresta $u\to v$ do grafo original passa a ser $u_{out}\to v_{in}$. Os caminhos mínimos no novo grafo correspondem aos do problema original. (Alternativa: some $c(v)$ ao peso de cada aresta que entra em $v$.)

\resp{dj-k} O algoritmo de Dijkstra só guarda o melhor custo até cada vértice, sem considerar quantas arestas foram usadas; o caminho mais barato até um vértice intermediário pode usar arestas demais. No grafo de estados, o vértice $(v,j)$ significa ``estou em $v$ depois de usar $j$ arestas'', com arestas $(u,j)\to(v,j+1)$ de peso $W(u,v)$ para $j<k$. A resposta é $\min_{j\le k}\operatorname{dist}((s,0),(t,j))$; como esse grafo é acíclico (o $j$ cresce), basta uma programação dinâmica em $O(k(n+m))$.

\resp{dj-gargalo} Troque $L(v)+W(v,w)$ por $\min\{L(v),\ cap(v,w)\}$, comece com $L(s)=\infty$ e $L(w)=0$ para os demais, e remova sempre o vértice de \emph{maior} rótulo. A prova de correção é análoga, porque estender um caminho nunca aumenta a sua capacidade (o papel dos pesos não negativos).

\resp{dj-conta} Mantenha $cnt(s)=1$ e, ao relaxar a aresta $vw$ ao remover $v$: se $L(v)+W(v,w)<L(w)$, faça $L(w):=L(v)+W(v,w)$ e $cnt(w):=cnt(v)$; se $L(v)+W(v,w)=L(w)$, faça $cnt(w):=cnt(w)+cnt(v)$. Com pesos positivos, quando $v$ é removido, $cnt(v)$ já está completo, pois todos os predecessores de $v$ em caminhos mínimos foram removidos antes dele.

\resp{dj-dag} Ordene os vértices topologicamente (Capítulo~\ref{cap:dfsdig}) e processe-os nessa ordem, relaxando as arestas que saem de cada um, a partir de $L(s)=0$. Quando um vértice é processado, todos os seus predecessores já foram, e o seu rótulo é definitivo. O custo é $O(n+m)$, e pesos negativos não são problema, pois não há ciclos.

\resp{dj-longo} A correção do algoritmo depende de que estender um caminho nunca o torne ``melhor''; para caminhos mais longos, estender sempre melhora, e um vértice fixado cedo pode ter um rótulo errado. Além disso, a melhor extensão pode precisar repetir vértices. O problema do caminho simples mais longo é NP-difícil em grafos gerais; só em grafos acíclicos ele é fácil.

\resp{dj-dial} Os rótulos finitos estão sempre entre $d$ e $d+C$, onde $d$ é o último rótulo removido. Use um vetor de baldes (listas) indexado pelo rótulo, de tamanho $nC+1$ (ou circular, de tamanho $C+1$). Para remover o mínimo, avance um ponteiro pelos baldes até encontrar um não vazio; o ponteiro só avança, no máximo $nC$ posições no total. Reduzir a chave é mover o vértice de um balde para outro, em $O(1)$. Total $O(m+nC)$.

\resp{dj-inv} Por indução: no início, $L(s)=0=\operatorname{dist}(s,s)$ e os demais valem $\infty$. Um relaxamento só muda $L(w)$ para $L(v)+W(v,w)\ge\operatorname{dist}(s,v)+W(v,w)\ge\operatorname{dist}(s,w)$, pela hipótese e pela desigualdade triangular; e só muda quando o novo valor é menor, de modo que $L(w)$ nunca aumenta.

\resp{dj-fib} No máximo $n$ inserções, $n$ remoções do mínimo e $m$ reduções de chave (no máximo uma por aresta, quando a aresta é relaxada). Com heap de Fibonacci, inserção e redução de chave custam $O(1)$ amortizado e remoção do mínimo $O(\log n)$ amortizado: o total é $O(n+m+n\log n)=O(m+n\log n)$.
